\documentclass[11pt]{article}

\usepackage[margin=1.1in]{geometry}
\usepackage{amsmath,amssymb}
\usepackage{booktabs}
\usepackage{enumitem}
\usepackage{xcolor}
\usepackage[colorlinks=true,linkcolor=black,urlcolor=blue,citecolor=black]{hyperref}

% --- status convention -------------------------------------------------
\newcommand{\status}[1]{\textcolor{gray}{\footnotesize\textsf{[#1]}}}
\newcommand{\verified}{\status{verified}}
\newcommand{\pending}{\status{pending fp32 rerun}}
\newcommand{\hypothesis}{\status{hypothesis, unevidenced}}
\newcommand{\open}{\status{open decision}}
\newcommand{\external}{\status{reported by DR, unverified by me}}
% -----------------------------------------------------------------------

\title{Static Fisher instruments on the Deep Ignorance checkpoints:\\
what is measured, what is claimed, and what is not}
\author{Alessa Carbo}
\date{Working note --- \today \\[2pt] \textsf{Draft. Numbers pending fp32 verification. Please do not cite.}}

\begin{document}
\maketitle

\begin{abstract}
\noindent This note states precisely which quantities I have measured on the
Deep Ignorance checkpoints, at what precision, under what scoping, and with
what open decisions outstanding. It is written to be self-contained: every
quantity is defined rather than referenced, so that the numbers can be checked
against a differently-implemented measurement without reconciling pipelines
first. Section~\ref{sec:discrepancy} addresses a specific disagreement between
my Fisher participation ratio measurements and Domenic's reported effective
Fisher rank on what may be the same checkpoint.
\end{abstract}

\paragraph{Status convention.} Every substantive claim below carries one of
five labels: \verified{} (measured, documented, reproducible from the notebook);
\pending{} (measured, but at a precision I am currently re-verifying);
\hypothesis{} (a mechanism I have proposed but not evidenced);
\open{} (a protocol choice I have not yet fixed); and \external{} (a claim
originating elsewhere that I am recording, not endorsing). Nothing in this note
is asserted at a higher confidence than its label.

\section{Definitions}

Since the point of this note is cross-checking, I define everything rather than
naming it.

\subsection{True vs.\ empirical Fisher}

Let $f_\theta(x)$ produce logits over an output space $\mathcal{Y}$, and write
$p_\theta(y \mid x)$ for the induced categorical distribution. The Fisher
information is
\begin{equation}
  I(\theta) \;=\; \mathbb{E}_{x \sim \mathcal{D}}\,
  \mathbb{E}_{y \sim p_\theta(\cdot \mid x)}
  \!\left[ \nabla_\theta \log p_\theta(y \mid x)\,
           \nabla_\theta \log p_\theta(y \mid x)^\top \right].
  \label{eq:true-fisher}
\end{equation}
The inner expectation is over labels drawn from the \emph{model's own}
predictive distribution. The \emph{empirical} Fisher replaces that inner
expectation with the observed label $y_i$:
\begin{equation}
  \tilde{I}(\theta) \;=\; \frac{1}{n}\sum_{i=1}^{n}
  \nabla_\theta \log p_\theta(y_i \mid x_i)\,
  \nabla_\theta \log p_\theta(y_i \mid x_i)^\top .
  \label{eq:emp-fisher}
\end{equation}
These coincide only near a well-specified optimum \citep{kunstner2019}. On a
forget distribution the model is \emph{deliberately} misspecified, which is
exactly the regime where the approximation is least defensible --- so the
distinction is load-bearing here rather than pedantic.

\paragraph{What I compute.} Equation~\eqref{eq:true-fisher}, exactly. The WMDP
battery is four-way multiple choice scored at the letter-logit tokens, so
$|\mathcal{Y}| = 4$ and the inner expectation is a sum over four terms rather
than a Monte Carlo estimate. There is no sampling error in the label
expectation. \verified{}

\subsection{Participation ratio}

Let $\lambda_1 \geq \dots \geq \lambda_d \geq 0$ be the eigenvalues of the
Fisher (or of the block being analysed). The participation ratio is
\begin{equation}
  \mathrm{PR} \;=\; \frac{\left(\sum_i \lambda_i\right)^2}{\sum_i \lambda_i^2}
  \;=\; \frac{(\operatorname{tr} I)^2}{\|I\|_F^2}.
  \label{eq:pr}
\end{equation}
It is scale-invariant, equals $d$ for a flat spectrum, and equals $1$ for a
rank-one spectrum. Reported here as $\log$-RMS across
\textsc{[placeholder: which layers, pooled or per-layer, and the exact
normalisation]}. \open{}

\subsection{Checkpoints}

\textsc{[placeholder: exact HuggingFace identifiers and revision hashes for the
three checkpoints --- strong- vs.\ weak-filter release, base vs.\
instruction-tuned. This is load-bearing given \S\ref{sec:threats}(3).]}

I refer to them as $M_{\mathrm{filt}}$ (filtered), $M_{\mathrm{know}}$
(intact), and $M_{\mathrm{CB}}$ (circuit breakers). Note the terminology
change: I previously called $M_{\mathrm{filt}}$ a ``never-knew'' baseline and
have stopped, for reasons in \S\ref{sec:threats}(3).

\subsection{Scope of the variance estimates}

Reported variances come from resampling forget questions on \emph{fixed}
checkpoints. They are \emph{not} across independently trained models. Every
equivalence claim below is therefore scoped to these specific checkpoints and
says nothing about the variability of the training procedure that produced
them. \verified{}

\section{Measurements}

\begin{table}[h]
\centering
\begin{tabular}{lccc}
\toprule
Checkpoint & $S_{\mathrm{all\ layers}}$ & $S_{\mathrm{excl.}}$ & Status \\
\midrule
$M_{\mathrm{filt}}$ & $0.339$ & $0.348$ & \pending{} \\
$M_{\mathrm{know}}$ & $0.332$ & $0.341$ & \pending{} \\
$M_{\mathrm{CB}}$   & $0.871$ & $0.738$ & \pending{}, \open{} (see \S\ref{sec:threats}(2)) \\
\bottomrule
\end{tabular}
\caption{Fisher participation-ratio summary across the three regimes. Both
variants use the conservative block-family set. $S_{\mathrm{all\ layers}}$
uses every eligible transformer layer; $S_{\mathrm{excl.}}$ excludes layers
30--31 for all three checkpoints. The CB separation survives both variants.}
\end{table}

\paragraph{What this supports.} The instrument separates the \emph{intervention
mechanism} --- $M_{\mathrm{CB}}$ stands clearly apart --- while failing to
separate having-known from never-having-known:
$M_{\mathrm{filt}} \approx M_{\mathrm{know}}$. Static curvature classifies
which intervention was applied, not whether a capability is present.
\pending{}

\paragraph{What this does not support.} Nothing about recoverability,
relearning speed, sample complexity to recovery, or attack outcomes. I have run
no optimization-axis measurement on these checkpoints. Any reading of these
numbers as evidence about tamper resistance is a misreading. \verified{}

\section{Known threats to the numbers}
\label{sec:threats}

\begin{enumerate}[leftmargin=*]

\item \textbf{Precision.} Current spectra were computed in
\textsc{[placeholder: precision of the existing harness]}. The concern is that
$\mathrm{PR}$ in~\eqref{eq:pr} is sensitive to the small-eigenvalue bulk: many
near-zero eigenvalues inflate $\operatorname{tr} I$ roughly linearly while
contributing negligibly to $\|I\|_F^2$, so $\mathrm{PR}$ rises. If that bulk
sits at or below a precision noise floor, part of what is being measured is the
floor --- and the floor is \emph{identical across checkpoints}, which is
precisely the shape needed to manufacture a spurious equivalence. An fp32 rerun
is scheduled. If $M_{\mathrm{filt}} \approx M_{\mathrm{know}}$ survives it, the
claim is strengthened; if it does not, the headline changes. \hypothesis{}

\item \textbf{CB layer-selection protocol.} Two protocols are in play:
per-cell maximum at \texttt{best\_layer}~24, and locked layers 30/26. They give
different gap magnitudes. Until one is fixed as canonical there is no single
quotable CB number. \open{}

\item \textbf{The reference model is contaminated.} $M_{\mathrm{filt}}$ is not
a never-knew model. Its pre-attack WMDP accuracy is already $\approx 0.42$
against a chance level of $0.25$, so the filter does not reach ignorance; and
Domenic reports rapid relearning in fp32 \external{}. Since the never-knew
baseline is the adjudicating anchor for the whole three-regime interpretation,
this matters: the ideal reference may not exist among realizable checkpoints,
which is a statement about the \emph{concept} and not only about this release.

\end{enumerate}

\section{A discrepancy worth resolving before either paper cites the other}
\label{sec:discrepancy}

Domenic reports that $M_{\mathrm{filt}}$ has \emph{collapsed effective Fisher
rank} (and minimal mutual information). I measure
$M_{\mathrm{filt}} \approx M_{\mathrm{know}}$ on Fisher-PR, i.e.\ not collapsed
relative to the intact model. These may or may not be compatible.

\paragraph{They are different functionals.} Effective rank is typically
\begin{equation}
  \mathrm{erank}(I) \;=\; \exp\!\big(H(\bar{\lambda})\big), \qquad
  \bar{\lambda}_i = \frac{\lambda_i}{\sum_j \lambda_j}, \qquad
  H(\bar\lambda) = -\sum_i \bar\lambda_i \log \bar\lambda_i,
\end{equation}
whereas $\mathrm{PR}$ is~\eqref{eq:pr}. Both are $1$ at rank one and $d$ at a
flat spectrum, but they weight the interior differently: $\mathrm{erank}$ is an
entropy and responds strongly to a long tail of small-but-nonzero eigenvalues,
while $\mathrm{PR}$ is an inverse second moment and is dominated by how much
mass sits in the top of the spectrum. Two spectra can therefore be ordered
oppositely by the two statistics. This is a genuine possibility and not an
excuse: if it resolves the disagreement, the fact that the choice of spectral
functional flips the story is itself a result the field should know.

\paragraph{Other candidate resolutions.}
\begin{enumerate}[label=(\alph*),leftmargin=*]
  \item different functional (above);
  \item true~\eqref{eq:true-fisher} vs.\ empirical~\eqref{eq:emp-fisher}
        Fisher --- the Kunstner critique applies asymmetrically here;
  \item precision, on either side;
  \item different evaluation distribution, layer set, or checkpoint variant.
\end{enumerate}

\paragraph{What I need from Domenic to close this.} The exact functional; true
or empirical Fisher; precision; evaluation distribution; layer set; and the
checkpoint identifier with revision. \paragraph{What I owe in return.} The same
list for my side (\S1), plus the fp32 rerun.

\section{What I have \emph{not} measured}

Recording this explicitly because a previous email of mine may have implied
otherwise.

\begin{itemize}[leftmargin=*]
  \item \textbf{RMU, RepNoise, TAR.} No Fisher measurements from me on these
        methods. My substrate is the three Deep Ignorance checkpoints; CB is
        the only overlap with the tamper-resistance method list. \verified{}
  \item \textbf{Mutual information.} Not run. Worth separating two objects: my
        enumeration gives \emph{exact} label-space $I(Y; \theta)$ over the
        four-outcome space with no estimator at all, which sidesteps the
        estimator pathologies entirely; it does \emph{not} give
        $I(\mathcal{D}_{\mathrm{forget}}; \theta)$, which is the quantity the
        certified-unlearning framing needs. Scoping question: does the
        multiclass MI you are using reduce to label-space MI? If so, exact
        values are cheap for me to produce.
  \item \textbf{Saddle index.} Not run. HVP infrastructure is in place, so the
        random-orthogonal-vector protocol should drop in once I have the spec.
\end{itemize}

\section{Provenance}

The numbers in \S2 currently trace to notebook runs. Figures lack committed
generating scripts, and the headline scalars are not persisted by a single
generator. Before these are used externally I will consolidate to one generator
emitting scalars, figures, and a SHA-256 input manifest, at fp32, against
pinned checkpoint revisions. This is the same work item as the fp32 rerun.
\open{}

\section*{Appendix, forthcoming separately}

A formalization of \emph{hidden} vs.\ \emph{removed}, drafted independently
per our pre-registration agreement: an optimization-axis definition, its
relativity to a reference class, the case where that reference class is empty,
and positioning against the certified-unlearning (DP-style) definition, which
is well-posed but presumes a realizable retrained reference.

\begin{thebibliography}{9}
\bibitem{kunstner2019}
F.~Kunstner, P.~Hennig, and L.~Balles.
\newblock Limitations of the empirical Fisher approximation for natural
gradient descent.
\newblock \emph{NeurIPS}, 2019.
\end{thebibliography}

\end{document}